% --- Archon physics formalization source begin ---
% archon:physics
% archon:covers PhyXMiniProblems/problem_phyx_mini_0212.lean
% archon:source-report reports/phyx_mini/problem_phyx_mini_0212.source.json

\chapter{Physics problem phyx\_mini\_0212}
\label{ch:PhyXMiniProblems_problem_phyx_mini_0212}

\paragraph{Problem source.}
\#\# Physical scenario

The sound level measured \textbackslash{}( 30\textbackslash{},\textbackslash{}mathrm\{m\} \textbackslash{}) from a jet plane is \textbackslash{}( 140\textbackslash{},\textbackslash{}mathrm\{dB\} \textbackslash{}).(Ignore reflections from the ground.)

\#\# Question

Estimate the sound level at \textbackslash{}( 300\textbackslash{},\textbackslash{}mathrm\{m\} \textbackslash{}).

\#\# Answer choices

- A: A: 160dB

- B: B: 140dB

- C: C: 100dB

- D: D: 120dB

\#\# Figure caption (auxiliary; use the image as primary evidence)

The image shows an airport ground crew worker directing an aircraft. The worker is wearing a high-visibility orange vest, navy blue sleeves with reflective stripes, and protective ear muffs. They are holding two-hand wands with red cones at the tips, which are used for signaling during aircraft operations. The person is standing in front of the airplane, which is partially visible, indicating a scene at an airport. The atmosphere is clear with a blue sky, and another airplane can be seen in the background to the right.

\#\# Dataset metadata

Sound; Physical Model Grounding Reasoning, Multi-Formula Reasoning

\paragraph{Recorded answer/context.}
D: D: 120dB

\paragraph{Figure/image path.}
/root/proposal\_for\_physic/science-mango/phyx\_mini\_run/phyx\_data/test\_image/212.png

\paragraph{Formalization target.}
create a compiling Lean file with sorry bodies at `PhyXMiniProblems/problem\_phyx\_mini\_0212.lean`. The Lean declarations must preserve the physical quantities, dimensions or dimensional roles, figure labels, governing-law hypotheses, and final relation expressed by this problem.
Use Mathlib/Physlib names found through LeanExplore where available. If a physics API is missing, introduce faithful local abstractions rather than scalar placeholder aliases.

\begin{theorem}[Physics formalization target]
\label{thm:physics:phyx_mini_0212:target}
\lean{PhyXMiniProblems.ProblemPhyXMini0212.problem_phyx_mini_0212}
\uses{def:physics:phyx-mini-0212:phyxminiproblems-problemphyxmini0212-jetplanesoundsetup, def:physics:phyx-mini-0212:phyxminiproblems-problemphyxmini0212-soundlevelindecibels, def:physics:phyx-mini-0212:phyxminiproblems-problemphyxmini0212-matchesjetplanescenario, def:physics:phyx-mini-0212:phyxminiproblems-problemphyxmini0212-matchessuppliedairportfigure, def:physics:phyx-mini-0212:phyxminiproblems-problemphyxmini0212-matchesproblemdata, def:physics:phyx-mini-0212:phyxminiproblems-problemphyxmini0212-usesstandardsoundintensityreference, def:physics:phyx-mini-0212:phyxminiproblems-problemphyxmini0212-hasphysicalacousticparameters, def:physics:phyx-mini-0212:phyxminiproblems-problemphyxmini0212-satisfiesisotropicinversesquarelaw, def:physics:phyx-mini-0212:phyxminiproblems-problemphyxmini0212-recordedanswerchoice, def:physics:phyx-mini-0212:phyxminiproblems-problemphyxmini0212-isclosestdisplayedanswer}
The assigned autoformalize agent should translate this physics problem into Lean declarations in the covered file, with theorem and lemma proof bodies written as `by sorry`.
\end{theorem}
\begin{proof}
This is an autoformalization task, not a proof task. Produce statements that can later be proved by the physics prover without weakening the physical model.
\end{proof}

% --- Archon declaration topology begin ---
\section{Declaration topology}
\begin{definition}[Length Magnitude]
  \label{def:physics:phyx-mini-0212:phyxminiproblems-problemphyxmini0212-lengthmagnitude}
  \lean{PhyXMiniProblems.ProblemPhyXMini0212.LengthMagnitude}
  A nonnegative physical source-to-observer distance
\end{definition}
\begin{proof}
  This is the declared model definition of Length Magnitude.
\end{proof}
\begin{definition}[Acoustic Power Quantity]
  \label{def:physics:phyx-mini-0212:phyxminiproblems-problemphyxmini0212-acousticpowerquantity}
  \lean{PhyXMiniProblems.ProblemPhyXMini0212.AcousticPowerQuantity}
  Acoustic power has SI unit watt and dimension mass * length\textasciicircum{}2 / time\textasciicircum{}3
\end{definition}
\begin{proof}
  This is the declared model definition of Acoustic Power Quantity.
\end{proof}
\begin{definition}[Acoustic Intensity Quantity]
  \label{def:physics:phyx-mini-0212:phyxminiproblems-problemphyxmini0212-acousticintensityquantity}
  \lean{PhyXMiniProblems.ProblemPhyXMini0212.AcousticIntensityQuantity}
  Acoustic intensity is power per area, with SI unit watt per square metre and dimension mass / time\textasciicircum{}3
\end{definition}
\begin{proof}
  This is the declared model definition of Acoustic Intensity Quantity.
\end{proof}
\begin{definition}[distance Readout]
  \label{def:physics:phyx-mini-0212:phyxminiproblems-problemphyxmini0212-distancereadout}
  \lean{PhyXMiniProblems.ProblemPhyXMini0212.distanceReadout}
  \uses{def:physics:phyx-mini-0212:phyxminiproblems-problemphyxmini0212-lengthmagnitude}
  Read a physical distance in a coherent choice of units
\end{definition}
\begin{proof}
  This is the declared model definition of distance Readout.
\end{proof}
\begin{definition}[acoustic Power Readout]
  \label{def:physics:phyx-mini-0212:phyxminiproblems-problemphyxmini0212-acousticpowerreadout}
  \lean{PhyXMiniProblems.ProblemPhyXMini0212.acousticPowerReadout}
  \uses{def:physics:phyx-mini-0212:phyxminiproblems-problemphyxmini0212-acousticpowerquantity}
  Read a physical acoustic power in a coherent choice of units
\end{definition}
\begin{proof}
  This is the declared model definition of acoustic Power Readout.
\end{proof}
\begin{definition}[acoustic Intensity Readout]
  \label{def:physics:phyx-mini-0212:phyxminiproblems-problemphyxmini0212-acousticintensityreadout}
  \lean{PhyXMiniProblems.ProblemPhyXMini0212.acousticIntensityReadout}
  \uses{def:physics:phyx-mini-0212:phyxminiproblems-problemphyxmini0212-acousticintensityquantity}
  Read a physical acoustic intensity in a coherent choice of units
\end{definition}
\begin{proof}
  This is the declared model definition of acoustic Intensity Readout.
\end{proof}
\begin{definition}[distance In Meters]
  \label{def:physics:phyx-mini-0212:phyxminiproblems-problemphyxmini0212-distanceinmeters}
  \lean{PhyXMiniProblems.ProblemPhyXMini0212.distanceInMeters}
  \uses{def:physics:phyx-mini-0212:phyxminiproblems-problemphyxmini0212-lengthmagnitude, def:physics:phyx-mini-0212:phyxminiproblems-problemphyxmini0212-distancereadout}
  SI metre readout of a source-to-observer distance
\end{definition}
\begin{proof}
  This is the declared model definition of distance In Meters.
\end{proof}
\begin{definition}[acoustic Power In Watts]
  \label{def:physics:phyx-mini-0212:phyxminiproblems-problemphyxmini0212-acousticpowerinwatts}
  \lean{PhyXMiniProblems.ProblemPhyXMini0212.acousticPowerInWatts}
  \uses{def:physics:phyx-mini-0212:phyxminiproblems-problemphyxmini0212-acousticpowerquantity, def:physics:phyx-mini-0212:phyxminiproblems-problemphyxmini0212-acousticpowerreadout}
  SI watt readout of an emitted acoustic power
\end{definition}
\begin{proof}
  This is the declared model definition of acoustic Power In Watts.
\end{proof}
\begin{definition}[acoustic Intensity In Watts Per Square Meter]
  \label{def:physics:phyx-mini-0212:phyxminiproblems-problemphyxmini0212-acousticintensityinwattspersquaremeter}
  \lean{PhyXMiniProblems.ProblemPhyXMini0212.acousticIntensityInWattsPerSquareMeter}
  \uses{def:physics:phyx-mini-0212:phyxminiproblems-problemphyxmini0212-acousticintensityquantity, def:physics:phyx-mini-0212:phyxminiproblems-problemphyxmini0212-acousticintensityreadout}
  SI watt-per-square-metre readout of an acoustic intensity
\end{definition}
\begin{proof}
  This is the declared model definition of acoustic Intensity In Watts Per Square Meter.
\end{proof}
\begin{definition}[Observation Point]
  \label{def:physics:phyx-mini-0212:phyxminiproblems-problemphyxmini0212-observationpoint}
  \lean{PhyXMiniProblems.ProblemPhyXMini0212.ObservationPoint}
  The measured and requested observer positions in the problem text
\end{definition}
\begin{proof}
  This is the declared model definition of Observation Point.
\end{proof}
\begin{definition}[Airport Figure Label]
  \label{def:physics:phyx-mini-0212:phyxminiproblems-problemphyxmini0212-airportfigurelabel}
  \lean{PhyXMiniProblems.ProblemPhyXMini0212.AirportFigureLabel}
  Qualitative objects visible in the supplied airport photograph
\end{definition}
\begin{proof}
  This is the declared model definition of Airport Figure Label.
\end{proof}
\begin{definition}[Airport Figure]
  \label{def:physics:phyx-mini-0212:phyxminiproblems-problemphyxmini0212-airportfigure}
  \lean{PhyXMiniProblems.ProblemPhyXMini0212.AirportFigure}
  \uses{def:physics:phyx-mini-0212:phyxminiproblems-problemphyxmini0212-airportfigurelabel}
  Visibility information extracted from the supplied airport photograph
\end{definition}
\begin{proof}
  This is the declared model definition of Airport Figure.
\end{proof}
\begin{definition}[Acoustic Source Kind]
  \label{def:physics:phyx-mini-0212:phyxminiproblems-problemphyxmini0212-acousticsourcekind}
  \lean{PhyXMiniProblems.ProblemPhyXMini0212.AcousticSourceKind}
  Type of idealized acoustic source used for the estimate
\end{definition}
\begin{proof}
  This is the declared model definition of Acoustic Source Kind.
\end{proof}
\begin{definition}[Emission Regime]
  \label{def:physics:phyx-mini-0212:phyxminiproblems-problemphyxmini0212-emissionregime}
  \lean{PhyXMiniProblems.ProblemPhyXMini0212.EmissionRegime}
  Time dependence of the jet's emitted acoustic power
\end{definition}
\begin{proof}
  This is the declared model definition of Emission Regime.
\end{proof}
\begin{definition}[Acoustic Medium]
  \label{def:physics:phyx-mini-0212:phyxminiproblems-problemphyxmini0212-acousticmedium}
  \lean{PhyXMiniProblems.ProblemPhyXMini0212.AcousticMedium}
  Propagation medium between the jet and the observers
\end{definition}
\begin{proof}
  This is the declared model definition of Acoustic Medium.
\end{proof}
\begin{definition}[Spreading Geometry]
  \label{def:physics:phyx-mini-0212:phyxminiproblems-problemphyxmini0212-spreadinggeometry}
  \lean{PhyXMiniProblems.ProblemPhyXMini0212.SpreadingGeometry}
  Geometric spreading model used for the sound field
\end{definition}
\begin{proof}
  This is the declared model definition of Spreading Geometry.
\end{proof}
\begin{definition}[Ground Reflection Treatment]
  \label{def:physics:phyx-mini-0212:phyxminiproblems-problemphyxmini0212-groundreflectiontreatment}
  \lean{PhyXMiniProblems.ProblemPhyXMini0212.GroundReflectionTreatment}
  How sound reflected from the ground is treated in the model
\end{definition}
\begin{proof}
  This is the declared model definition of Ground Reflection Treatment.
\end{proof}
\begin{definition}[Jet Plane Sound Setup]
  \label{def:physics:phyx-mini-0212:phyxminiproblems-problemphyxmini0212-jetplanesoundsetup}
  \lean{PhyXMiniProblems.ProblemPhyXMini0212.JetPlaneSoundSetup}
  \uses{def:physics:phyx-mini-0212:phyxminiproblems-problemphyxmini0212-lengthmagnitude, def:physics:phyx-mini-0212:phyxminiproblems-problemphyxmini0212-acousticpowerquantity, def:physics:phyx-mini-0212:phyxminiproblems-problemphyxmini0212-acousticintensityquantity, def:physics:phyx-mini-0212:phyxminiproblems-problemphyxmini0212-observationpoint, def:physics:phyx-mini-0212:phyxminiproblems-problemphyxmini0212-airportfigure, def:physics:phyx-mini-0212:phyxminiproblems-problemphyxmini0212-acousticsourcekind, def:physics:phyx-mini-0212:phyxminiproblems-problemphyxmini0212-emissionregime, def:physics:phyx-mini-0212:phyxminiproblems-problemphyxmini0212-acousticmedium, def:physics:phyx-mini-0212:phyxminiproblems-problemphyxmini0212-spreadinggeometry, def:physics:phyx-mini-0212:phyxminiproblems-problemphyxmini0212-groundreflectiontreatment}
  Physical quantities for the two-position jet-sound estimate. The intensity at the far observation point is an independent physical field; neither its numerical value nor the requested far sound level is stored in this structure
\end{definition}
\begin{proof}
  This is the declared model definition of Jet Plane Sound Setup.
\end{proof}
\begin{definition}[sound Level In Decibels]
  \label{def:physics:phyx-mini-0212:phyxminiproblems-problemphyxmini0212-soundlevelindecibels}
  \lean{PhyXMiniProblems.ProblemPhyXMini0212.soundLevelInDecibels}
  \uses{def:physics:phyx-mini-0212:phyxminiproblems-problemphyxmini0212-acousticintensityinwattspersquaremeter, def:physics:phyx-mini-0212:phyxminiproblems-problemphyxmini0212-observationpoint, def:physics:phyx-mini-0212:phyxminiproblems-problemphyxmini0212-jetplanesoundsetup}
  The sound-intensity level at an observation point, in decibels. This is the standard generic definition 10 log10 (I / I\_ref) applied to same-unit SI readouts; it contains no distance-specific numerical answer
\end{definition}
\begin{proof}
  This is the declared model definition of sound Level In Decibels.
\end{proof}
\begin{definition}[Matches Jet Plane Scenario]
  \label{def:physics:phyx-mini-0212:phyxminiproblems-problemphyxmini0212-matchesjetplanescenario}
  \lean{PhyXMiniProblems.ProblemPhyXMini0212.MatchesJetPlaneScenario}
  \uses{def:physics:phyx-mini-0212:phyxminiproblems-problemphyxmini0212-jetplanesoundsetup}
  The idealized physical scenario used for the estimate: a constant-power jet point source radiates spherically through ambient air, and ground reflections are omitted as required by the problem statement
\end{definition}
\begin{proof}
  This is the declared model definition of Matches Jet Plane Scenario.
\end{proof}
\begin{definition}[Matches Supplied Airport Figure]
  \label{def:physics:phyx-mini-0212:phyxminiproblems-problemphyxmini0212-matchessuppliedairportfigure}
  \lean{PhyXMiniProblems.ProblemPhyXMini0212.MatchesSuppliedAirportFigure}
  \uses{def:physics:phyx-mini-0212:phyxminiproblems-problemphyxmini0212-jetplanesoundsetup}
  Qualitative evidence in the primary image: a ground-crew worker wearing hearing protection raises two signal wands beside a foreground aircraft, with another aircraft in the background. The image supplies no numerical distance or sound-level readout
\end{definition}
\begin{proof}
  This is the declared model definition of Matches Supplied Airport Figure.
\end{proof}
\begin{definition}[Matches Problem Data]
  \label{def:physics:phyx-mini-0212:phyxminiproblems-problemphyxmini0212-matchesproblemdata}
  \lean{PhyXMiniProblems.ProblemPhyXMini0212.MatchesProblemData}
  \uses{def:physics:phyx-mini-0212:phyxminiproblems-problemphyxmini0212-distanceinmeters, def:physics:phyx-mini-0212:phyxminiproblems-problemphyxmini0212-jetplanesoundsetup, def:physics:phyx-mini-0212:phyxminiproblems-problemphyxmini0212-soundlevelindecibels}
  The numerical data stated in the problem. The near measurement is 140 dB at 30 m, while the requested observation position is 300 m away. No far sound level or answer choice occurs in these data
\end{definition}
\begin{proof}
  This is the declared model definition of Matches Problem Data.
\end{proof}
\begin{definition}[Uses Standard Sound Intensity Reference]
  \label{def:physics:phyx-mini-0212:phyxminiproblems-problemphyxmini0212-usesstandardsoundintensityreference}
  \lean{PhyXMiniProblems.ProblemPhyXMini0212.UsesStandardSoundIntensityReference}
  \uses{def:physics:phyx-mini-0212:phyxminiproblems-problemphyxmini0212-acousticintensityinwattspersquaremeter, def:physics:phyx-mini-0212:phyxminiproblems-problemphyxmini0212-jetplanesoundsetup}
  The conventional reference intensity for sound intensity level is 10\textasciicircum{}-12 W/m\textasciicircum{}2. This convention is independent of either observation distance and does not determine the requested far level by itself
\end{definition}
\begin{proof}
  This is the declared model definition of Uses Standard Sound Intensity Reference.
\end{proof}
\begin{definition}[Has Physical Acoustic Parameters]
  \label{def:physics:phyx-mini-0212:phyxminiproblems-problemphyxmini0212-hasphysicalacousticparameters}
  \lean{PhyXMiniProblems.ProblemPhyXMini0212.HasPhysicalAcousticParameters}
  \uses{def:physics:phyx-mini-0212:phyxminiproblems-problemphyxmini0212-distanceinmeters, def:physics:phyx-mini-0212:phyxminiproblems-problemphyxmini0212-acousticpowerinwatts, def:physics:phyx-mini-0212:phyxminiproblems-problemphyxmini0212-acousticintensityinwattspersquaremeter, def:physics:phyx-mini-0212:phyxminiproblems-problemphyxmini0212-observationpoint, def:physics:phyx-mini-0212:phyxminiproblems-problemphyxmini0212-jetplanesoundsetup}
  Strict positivity and nondegeneracy of the physical acoustic quantities
\end{definition}
\begin{proof}
  This is the declared model definition of Has Physical Acoustic Parameters.
\end{proof}
\begin{definition}[Satisfies Isotropic Inverse Square Law]
  \label{def:physics:phyx-mini-0212:phyxminiproblems-problemphyxmini0212-satisfiesisotropicinversesquarelaw}
  \lean{PhyXMiniProblems.ProblemPhyXMini0212.SatisfiesIsotropicInverseSquareLaw}
  \uses{def:physics:phyx-mini-0212:phyxminiproblems-problemphyxmini0212-distancereadout, def:physics:phyx-mini-0212:phyxminiproblems-problemphyxmini0212-acousticpowerreadout, def:physics:phyx-mini-0212:phyxminiproblems-problemphyxmini0212-acousticintensityreadout, def:physics:phyx-mini-0212:phyxminiproblems-problemphyxmini0212-observationpoint, def:physics:phyx-mini-0212:phyxminiproblems-problemphyxmini0212-jetplanesoundsetup}
  The free-field inverse-square intensity law for an isotropic point source. In every coherent unit system, the acoustic power crossing the sphere of radius r is 4 * pi * r\textasciicircum{}2 * I. This general law contains neither the 1/100 far-intensity result nor the requested 120 dB conclusion
\end{definition}
\begin{proof}
  This is the declared model definition of Satisfies Isotropic Inverse Square Law.
\end{proof}
\begin{lemma}[far Intensity eq near Intensity div hundred]
  \label{lem:physics:phyx-mini-0212:phyxminiproblems-problemphyxmini0212-farintensity-eq-nearintensity-div-hundred}
  \lean{PhyXMiniProblems.ProblemPhyXMini0212.farIntensity_eq_nearIntensity_div_hundred}
  \uses{def:physics:phyx-mini-0212:phyxminiproblems-problemphyxmini0212-acousticintensityinwattspersquaremeter, def:physics:phyx-mini-0212:phyxminiproblems-problemphyxmini0212-jetplanesoundsetup, def:physics:phyx-mini-0212:phyxminiproblems-problemphyxmini0212-matchesproblemdata, def:physics:phyx-mini-0212:phyxminiproblems-problemphyxmini0212-hasphysicalacousticparameters, def:physics:phyx-mini-0212:phyxminiproblems-problemphyxmini0212-satisfiesisotropicinversesquarelaw}
  Increasing the distance from 30 m to 300 m reduces the acoustic intensity by a factor of 100. This is a derived result, not a model assumption
\end{lemma}
\begin{proof}
  The conclusion follows from the governing relations and figure readouts named in the statement of far Intensity eq near Intensity div hundred.
\end{proof}
\begin{lemma}[far Sound Level eq near Sound Level sub twenty]
  \label{lem:physics:phyx-mini-0212:phyxminiproblems-problemphyxmini0212-farsoundlevel-eq-nearsoundlevel-sub-twenty}
  \lean{PhyXMiniProblems.ProblemPhyXMini0212.farSoundLevel_eq_nearSoundLevel_sub_twenty}
  \uses{def:physics:phyx-mini-0212:phyxminiproblems-problemphyxmini0212-jetplanesoundsetup, def:physics:phyx-mini-0212:phyxminiproblems-problemphyxmini0212-soundlevelindecibels, def:physics:phyx-mini-0212:phyxminiproblems-problemphyxmini0212-matchesproblemdata, def:physics:phyx-mini-0212:phyxminiproblems-problemphyxmini0212-hasphysicalacousticparameters, def:physics:phyx-mini-0212:phyxminiproblems-problemphyxmini0212-satisfiesisotropicinversesquarelaw}
  The inverse-square intensity change across a tenfold distance increase lowers the sound intensity level by 20 dB. This conclusion is derived directly from the problem data and governing law
\end{lemma}
\begin{proof}
  The conclusion follows from the governing relations and figure readouts named in the statement of far Sound Level eq near Sound Level sub twenty.
\end{proof}
\begin{definition}[Answer Choice]
  \label{def:physics:phyx-mini-0212:phyxminiproblems-problemphyxmini0212-answerchoice}
  \lean{PhyXMiniProblems.ProblemPhyXMini0212.AnswerChoice}
  Labels of the four displayed multiple-choice answers
\end{definition}
\begin{proof}
  This is the declared model definition of Answer Choice.
\end{proof}
\begin{definition}[decibels]
  \label{def:physics:phyx-mini-0212:phyxminiproblems-problemphyxmini0212-answerchoice-decibels}
  \lean{PhyXMiniProblems.ProblemPhyXMini0212.AnswerChoice.decibels}
  \uses{def:physics:phyx-mini-0212:phyxminiproblems-problemphyxmini0212-answerchoice}
  Sound-level readout in decibels printed beside each answer label
\end{definition}
\begin{proof}
  This is the declared model definition of decibels.
\end{proof}
\begin{definition}[recorded Answer Choice]
  \label{def:physics:phyx-mini-0212:phyxminiproblems-problemphyxmini0212-recordedanswerchoice}
  \lean{PhyXMiniProblems.ProblemPhyXMini0212.recordedAnswerChoice}
  \uses{def:physics:phyx-mini-0212:phyxminiproblems-problemphyxmini0212-answerchoice}
  The answer label recorded by the supplied dataset
\end{definition}
\begin{proof}
  This is the declared model definition of recorded Answer Choice.
\end{proof}
\begin{definition}[Is Closest Displayed Answer]
  \label{def:physics:phyx-mini-0212:phyxminiproblems-problemphyxmini0212-isclosestdisplayedanswer}
  \lean{PhyXMiniProblems.ProblemPhyXMini0212.IsClosestDisplayedAnswer}
  \uses{def:physics:phyx-mini-0212:phyxminiproblems-problemphyxmini0212-answerchoice}
  A displayed choice is uniquely closest to a computed sound-level estimate. This generic comparison does not privilege the recorded choice by definition
\end{definition}
\begin{proof}
  This is the declared model definition of Is Closest Displayed Answer.
\end{proof}
% --- Archon declaration topology end ---
% --- Archon physics formalization source end ---
